\section{The natural lattice and the depth of an opposite extension}\label{sec:depth}
The reducible branch uses two different kinds of radical. The base field is
\[
K=\Q(\sqrt5),\qquad\phi=(1+\sqrt5)/2.
\]
The coefficient ring and field are
\[
\Oloc=\Z_7[w]/(w^2-5),\qquad k=\F_7[w]/(w^2-5)=\F_{49}.
\]
The base unit $\phi$ defines Kummer classes. The coefficient $w$ acts on a rank two representation. Although both involve a square root of five, they have different roles and are never identified by an implicit residue substitution.

The inherited characteristic zero representation is the specified RM factor of the twisted Jacobian on the retained $3990\mid C$ branch, with determinant the cyclotomic character $\chi$. Its natural residual orientation is
\[
0\longrightarrow k(\chi)\longrightarrow\overline V\longrightarrow k\longrightarrow0.
\]
The opposite endpoint has a nonzero class $\beta\in H^1(K,k(\chi^{-1}))$. A stable lattice change can alter residual inertia. Therefore properties of the natural Jacobian lattice cannot automatically be transferred to an arbitrary lattice supplied by Ribet's lemma \cite{ophir-weiss2024}.

\subsection{Why a new eligible prime is necessary}
The opposite orientation forces an additional prime divisor of $C$. Here is the conditional argument, retaining the class group and local geometric inputs. Suppose a characteristic zero local toric line has reduction $\chi$, while the opposite residual extension has subline one and quotient $\chi$. Intersect the actual toric line with its stable lattice. The intersection is saturated, so it reduces to a nonzero local $\chi$ line. If the local character $\chi$ is nontrivial, that line is distinct from the global subline one. They split the residual extension locally. This proof survives stable lattice changes because it starts with the characteristic zero line.

At seven, the specified twist makes the toric characters $\chi$ and one, and cyclotomic inertia is nontrivial. At a variable bad place with $Q\not\equiv1\pmod7$, the same argument applies. At the places over two, three and five, a separate local cohomology calculation excludes opposite extensions. For a non seven local field of residue cardinality $Q$, an inertia cocycle with coefficients $k(\chi^{-1})$ can be nonzero only when $Q^2\equiv1\pmod7$: the tame Frobenius relation multiplies its inertia value by $Q$, while the coefficient action multiplies it by $Q^{-1}$. The unramified quotient is zero when $Q\not\equiv1$. Thus $Q\not\equiv\pm1$ forces local $H^1$ to vanish. Here the relevant norms are four, nine and five, all outside those two residue classes.

If no variable $C$ prime had $Q\equiv1\pmod7$, beta would restrict to an everywhere unramified additive class over $L=K(\zeta_7)$. Restriction is injective because $[L:K]=6$ is prime to seven. The inherited Masley class number result gives $h(\Q(\zeta_{35}))=1$ \cite{masley1976}. The field $L$ is a subfield of index two. Extension and norm on ideal classes compose to multiplication by two, so the seven primary class group of $L$ is zero. An everywhere unramified nonzero additive seven class is therefore impossible by class field theory.

At least one variable prime above a divisor of $C$ must consequently have norm one modulo seven. A rational prime splits in $K$ when it is one or four modulo five, giving residue degree one. At an inert prime the degree is two. Combining $q^f\equiv1\pmod7$ with these splitting possibilities gives exactly
\[
q\pmod{35}\in\{1,8,13,22,27,29\}.
\]
None of $2,3,5,7,19$, the primes already in $3990$, supplies this new variable prime. The conclusion asks for an eligible additional prime, not divisibility by a particular fixed prime. Once such primes are allowed, their ray class columns can support beta. Ordinary class number one does not make that ramified character space vanish.

\subsection{An actual finite segment of stable lattices}
The inherited split toric calculation at three supplies Frobenius eigenvalues nine and one. Their difference eight is a seven adic unit. The projectors
\[
E_9=\frac{\rho(g_0)-I}{8},\qquad E_1=\frac{9I-\rho(g_0)}8
\]
are complementary integral idempotents. Choose generators of their images. In this basis, $\rho(g_0)=\operatorname{diag}(9,1)$ and every lower entry is divisible by seven. Natural upper monodromy at three has a unit entry. Absolute irreducibility in characteristic zero forces some lower entry to be nonzero.

Define
\begin{equation}\label{eq:actualdepth}
n=\min_{g:\rho(g)_{21}\ne0}v_7(\rho(g)_{21}).
\end{equation}
This positive integer is attained. It is an invariant of the fixed natural frame, not a truncation selected for a computation. Every $\Oloc$ lattice stable under $\rho$ is preserved by the two projectors, so up to homothety it has the form
\[
\Lambda_j=\Oloc e_1+\Oloc 7^j e_2.
\]
The unit upper monodromy forces $j\ge0$. The attained lower entry minimum forces $j\le n$. Conversely every $j$ in that interval is stable. Thus the stable lattices form a finite segment with the natural and opposite endpoints at $0,n$.

At the opposite endpoint, the attained minimum gives a nonzero lower residual entry. A split complement would have to be the first coordinate line, because $g_0$ has distinct residual eigenvalues. That line is not stable. Hence the opposite residual extension is nonsplit. This constructs the relevant class on a compatible endpoint rather than importing natural inertia through an unjustified lattice change.

\subsection{A local depth theorem}
Here is the matrix argument controlling where this opposite class can ramify. Let $\Oloc$ be a discrete valuation ring with uniformizer $\pi$, let $F\in\mathrm{GL}_2(\Oloc)$, $Q\in\Oloc^\times$ with $Q\ne1$, and let $n\ge1$. Suppose $N\in M_2(\Oloc)$ is a nonzero nilpotent matrix of the form
\[
N=\mu\begin{pmatrix}-xy&x^2\\-y^2&xy\end{pmatrix},
\]
where $\mu\in\Oloc\setminus\{0\}$, $(x,y)\in\Oloc^2$ is primitive, $t=v(\mu)\ge0$, and
\[
FN=QN,\qquad NF=N,\qquad F_{21}\in\pi^n\Oloc,\qquad v(N_{21})=n.
\]
Set $h=v(Q-1)$. Then
\begin{equation}\label{eq:depthlocal}
n=t+2s,\qquad s=v(y),\qquad0\le s\le h.
\end{equation}
\begin{proof}
The lower nilpotent entry is $-\mu y^2$, so its valuation gives $n=t+2s$. The common image and kernel line of $N$ is generated by $(x,y)^T$. The two identities for $F$ show that $F-I$ has image in this line and kills no extra quotient direction. Write
\[
F=I+\begin{pmatrix}x\\y\end{pmatrix}(a,b),\qquad ax+by=Q-1.
\]
The lower entry is $ya$, so $a\in\pi^{n-s}\Oloc=\pi^{t+s}\Oloc$. Both $ax$ and $by$ are divisible by $\pi^s$. Their sum $Q-1$ is therefore divisible by $\pi^s$, giving $s\le h$.
\end{proof}
The square on $y$ explains the factor two. A bound $n\le t+h$ would omit the geometry of the rank one nilpotent and is false. The examples
\[
N=7^t\begin{pmatrix}-7^s&1\\-7^{2s}&7^s\end{pmatrix},\qquad
F=\begin{pmatrix}1&(Q-1)/7^s\\0&Q\end{pmatrix}
\]
realize every allowed $s$ when $v_7(Q-1)=h$. They are integral local matrix controls. They are not automatically global Galois representations or Jacobians.

For the natural curve, the variable node thickness is $7v_q(C)$. The natural monodromy content is consequently
\[
t=1+v_7(v_q(C)).
\]
This inherited geometric bridge uses the ordinary node smoothing parameter and the identification of the graph pairing with the period valuation \cite{baker-rabinoff2013}. Applied to an actually ramified beta prime, \eqref{eq:depthlocal} becomes
\begin{equation}\label{eq:arithdepth}
n=1+v_7(v_q(C))+2s,\qquad0\le s\le v_7(q^f-1).
\end{equation}
Here $f$ is the residue degree in $K$. If $7\nmid v_q(C)$ at every eligible prime, the actual depth is odd. If in addition every eligible residue norm is nontrivial modulo forty nine, the only possible depths are one and three. These are necessary conditions within the inherited arithmetic contract, not a uniform depth bound for arbitrary $C$.

\subsection{The local cup endpoint}
At an eligible beta prime, $t\ge1$ and $h\ge1$, so $\chi$ is trivial modulo seven. Natural inertia is trivial modulo seven, while opposite inertia is nonzero. Below the upper local endpoint $n=t+2h$, the same matrix calculation forces natural Frobenius $\alpha(F)$ to vanish. At the endpoint it is nonzero.

If $\alpha\cup\beta$ were a coboundary in the endpoint case, the two adjacent entries would extend to a representation in the three by three upper unitriangular group over $k$. That group has exponent seven. The tame relation $FTF^{-1}=T^Q$ with $Q\equiv1\pmod7$ would force the images to commute. Their central commutator is $\alpha(F)\beta(T)-\alpha(T)\beta(F)=\alpha(F)\beta(T)\ne0$. This contradiction proves local nonvanishing at the endpoint. The zero below the endpoint and nonzero at it are statements about the same labelled local representation.
